\section{Applications}
Thanks to the powerful representation of superpoints for dynamic scenes, our method is highly expandable and can facilitate various downstream applications. 
\subsection{Model Distillation}
In scenarios where a 3D-GS based model $\mathcal{H}$ exhibits superior performance, which also predicts the deformation over time, we can distill such a model into SP-GS to improve the visual quality.
The concept is straightforward: we can directly replicate the state of $\mathcal{H}$ at any given time as the canonical state of SP-GS. 
Subsequently, we optimize the association matrix $\Asso$, the superpoint deformation network $\mathcal{F}$, and optionally the non-rigid deformation network $\mathcal{G}$ by incorporating $\loss$ loss and mean square error(MSE) $ \loss_{err} =\sum_{\bm{u} \in \{\gsP, \gsR \}} \lambda_{\bm{u}} \sum_{i} \|\bm{u}^t_i, \bm{u}'^t_i \|$,
% \begin{equation}
%     \loss_{err} =\sum_{\bm{u} \in \{\gsP, \gsR \}} \lambda_{\bm{u}} \sum_{i} \|\bm{u}^t_i, \bm{u}'^t_i \|,
% \end{equation}
where $\bm{u}^t_i$  and $\bm{u}'^t_i$ are the properties of the teacher and the student respectively.
Tab. \ref{tab:distill} shows the quantitative results of distilling D-3D-GS model into SP-GS model on the ``As" scene of NeRF-DS dataset.
While D-3D-GS cannot achieve real-time rendering on V100 (20.65 FPS), our distillated student model can achieve significantly higher rendering speed (164.04 FPS). Therefore, model distillation provides a trade-off between visual quality and rendering speed, leaving users with more choices to meet their requirements.

\subsection{Pose Estimation}

Our SP-GS support estimating the 6-DoF pose of each superpoint for new images in the same scene.
%  without relying on the deformation network. 
To be specific, we can solely learn the translation and rotation for each superpoint with other parameters of SP-GS fixed. This can be potentially used in motion capture where novel view images are given and one wants to know the motion of each components (superpoints).
Experiments are conducted on the jumpingjacks scene of D-NeRF dataset~\cite{D-NeRF}. We first train the complete model (SP-GS) using the beginning 50 images. Subsequently, we initialize the learnable translation and rotation parameters of superpoints with the 50-th frame and directly optimize them with rendering loss using the Adam optimizer for 1000 iterations.
The program terminates upon completing the pose estimation for all images.
Fig. \ref{fig:repose} illustrates the change in PSNR across images 51-88 for the jumpingjacks scene.
Notably, it demonstrates a gradual decrease in PSNR.

\subsection{Scene Editing}
As depicted in Figure \ref{fig:edit}, scene editing tasks such as relocating parts between scenes or removing parts from a scene can be accomplished with ease. 
This capability is facilitated by the explicit 3D Gaussian representation, enabling relocation or deletion from the scene.
Our method further streamline the process since it is no longer necessary to manipulate over the 100,000 3D Gaussians. Moreover, our superpoints provide some extent of semantic meanings, enabling a reasonable editing of the scene.
